% ==============================================================================
% MANUAL OPERATIVO: TECNICAS DIDACTICAS DEL APRENDIZAJE
% Fuente de referencia: Coleccion "100 Tecnicas Didacticas de Ensenanza y
% Aprendizaje", Fasciculos 1 a 5, UnADM, 2023.
%
% Proposito del archivo:
% - Servir como guia para elegir, construir y revisar productos didacticos.
% - Traducir la tecnica solicitada en un producto LaTeX/TikZ verificable.
% - Evitar entregas genericas: todo producto debe tener estructura, criterio,
%   analisis y postura cuando la actividad lo permita.
%
% Regla general para nuevas actividades:
% Si la planeacion pide tabla, cuadro comparativo, mapa conceptual, mapa mental,
% linea de tiempo, diagrama, esquema, matriz, cronologia, red conceptual,
% histograma, flujo o proceso, construirlo preferentemente con TikZ/LaTeX.
% Usar captura externa solo si la consigna exige una herramienta especifica o si
% el producto requiere una plataforma interactiva.
% ==============================================================================

\clearpage
\section{Manual operativo de técnicas didácticas}

\subsection{Regla base para productos visuales}

Cuando la actividad solicite un producto visual, el producto debe construirse como parte del documento y no como un elemento improvisado. La regla de trabajo es: si el producto puede representarse con nodos, flechas, tablas, matrices, líneas, ciclos, mapas o relaciones, debe elaborarse en TikZ o con entornos LaTeX como \texttt{tabular}, \texttt{longtable} o \texttt{figure}. Esto permite controlar formato, legibilidad, jerarquía, citas, continuidad visual y coherencia con el documento.

El producto no debe funcionar como adorno. Debe demostrar comprensión: seleccionar información, jerarquizarla, relacionarla y cerrar con una lectura propia. En términos prácticos, cada producto debe contestar cuatro preguntas: qué organiza, con qué criterios lo organiza, qué relación muestra y qué conclusión permite sostener.

\subsection{Estructura común de las técnicas UnADM}

Los fascículos de \textit{100 Técnicas Didácticas de Enseñanza y Aprendizaje} trabajan cada técnica con una lógica estable: definición, estructura, utilidad, proceso de construcción, consideraciones y referencias. Para convertir esa lógica en una actividad universitaria, se recomienda usar esta secuencia:

\begin{enumerate}
	\item \textbf{Identificar el producto solicitado.} No es lo mismo comparar, explicar, analizar, sintetizar o construir. El verbo de la consigna define el nivel de profundidad.
	\item \textbf{Determinar los elementos obligatorios.} Toda técnica tiene piezas mínimas: categorías, conceptos, fuentes, etapas, actores, criterios, datos, ejemplos o conclusiones.
	\item \textbf{Elegir el formato TikZ/LaTeX.} Tabla para comparar, nodos para mapas, flechas para procesos, línea horizontal para cronologías, matriz para decisiones.
	\item \textbf{Redactar antes de diseñar.} El diseño no corrige una idea débil. Primero se definen conceptos y relaciones; después se dibuja.
	\item \textbf{Agregar lectura crítica.} El producto debe mostrar qué significa la información, por qué importa y cómo se aplica.
	\item \textbf{Cerrar con postura.} Si la actividad pide conclusión, debe responder de forma directa la pregunta final y justificarla.
\end{enumerate}

\subsection{Taxonomía de ejecución}

La colección clasifica las técnicas por nivel de ejecución. Esta lectura ayuda a decidir qué profundidad debe tener la entrega:

\begin{longtable}{@{}p{0.14\linewidth}p{0.2\linewidth}p{0.56\linewidth}@{}}
\toprule
\textbf{Nivel} & \textbf{Acción} & \textbf{Criterio para la actividad} \\
\midrule
\endhead
1 & Recordar & Recuperar datos, conceptos, nombres o hechos. Requiere precisión, pero no basta para una actividad universitaria completa. \\
2 & Explicar & Aclarar el significado de un concepto y su función. Debe incluir ejemplos o aplicación. \\
3 & Aplicar & Usar el conocimiento en una situación, caso, norma o problema. Debe mostrar procedimiento. \\
4 & Analizar & Separar elementos, identificar causas, efectos, relaciones y límites. Exige postura crítica. \\
5 & Sintetizar & Integrar información dispersa en un producto ordenado. Debe mostrar jerarquía y selección. \\
6 & Construir & Crear un producto, propuesta, estrategia, guion, ensayo, proyecto o solución. Exige diseño completo y justificación. \\
\bottomrule
\end{longtable}

\subsection{Manual por producto frecuente}

\subsubsection*{Cuadro comparativo}

Usar cuando la consigna pida comparar corrientes, autores, normas, escuelas, conceptos, textos, versiones o modelos. Debe contener una tabla con conceptos en columnas y categorías en filas. Las categorías no deben ser genéricas: deben permitir contraste real. Para asignaturas de Derecho se recomiendan categorías como concepto, fundamento, autores, principios, función, aplicación, límites y valoración crítica; para Redacción en contextos virtuales conviene usar claridad, coherencia, cohesión, formalidad, objetividad, adecuación al receptor y efecto comunicativo.

\textbf{Construcción:} investigar, seleccionar información relevante, definir categorías, ordenar por jerarquía, redactar celdas breves pero sustanciales, citar fuentes y cerrar con conclusión personal. En TikZ/LaTeX puede hacerse con \texttt{longtable}; si se requiere diseño más visual, usar nodos en matriz.

\textbf{Revisión:} cada fila debe permitir una comparación inmediata; si una categoría no cambia la comprensión, se elimina o se reemplaza.

\subsubsection*{Cuadro sinóptico}

Sirve para organizar un tema de lo general a lo particular. Debe mostrar niveles: tema central, categorías principales, subcategorías y detalles. Es útil para corrientes filosóficas, clasificación de derechos, teorías de justicia o tipos de normas.

\textbf{Construcción:} definir tema central, extraer categorías, ordenar por dependencia lógica y usar llaves, ramas o cajas. En TikZ se recomienda una estructura de árbol con nodos alineados.

\subsubsection*{Mapa conceptual}

Debe mostrar conceptos y relaciones mediante palabras de enlace. No basta con poner cajas conectadas; cada flecha debe decir qué relación existe: causa, implica, deriva, limita, permite, exige, se opone a, fundamenta.

\textbf{Construcción:} seleccionar concepto central, jerarquizar conceptos principales y secundarios, escribir conectores precisos, evitar saturación visual y cerrar con explicación breve del mapa. En TikZ usar nodos jerárquicos con flechas etiquetadas.

\subsubsection*{Mapa mental}

Funciona para explorar asociaciones desde una idea central. A diferencia del mapa conceptual, admite ramas más libres, colores, imágenes o palabras clave. Es útil para lluvia de ideas, conceptos iniciales o planeación de investigación.

\textbf{Construcción:} ubicar idea central, crear ramas principales, agregar subramas, usar palabras breves y cerrar con una interpretación. En TikZ usar nodos radiales o coordenadas polares.

\subsubsection*{Línea de tiempo y cronología ilustrada}

La línea de tiempo organiza eventos por secuencia temporal; la cronología ilustrada agrega apoyo visual o iconográfico. Sirve para evolución del pensamiento jurídico, historia de derechos, reformas constitucionales o desarrollo de una corriente.

\textbf{Construcción:} seleccionar fechas relevantes, ordenar cronológicamente, resumir cada evento, señalar causa o consecuencia y cerrar con lectura histórica. En TikZ usar una línea horizontal con hitos alternados arriba/abajo.

\subsubsection*{Esquema}

El esquema resume una estructura lógica. Puede ser jerárquico, secuencial, causal o temático. Sirve cuando la actividad pide explicar un tema sin llegar a mapa conceptual.

\textbf{Construcción:} identificar tema, dividirlo en partes, ordenar relaciones y redactar con frases breves. En TikZ usar cajas alineadas o árbol simple.

\subsubsection*{Diagrama de flujo}

Representa pasos, decisiones y resultados. Es útil para procedimientos jurídicos, interpretación normativa, juicio de amparo, análisis de caso o toma de decisiones.

\textbf{Construcción:} definir inicio, pasos, decisiones, salidas y cierre. Usar rombos para decisiones, rectángulos para acciones y flechas para secuencia. En TikZ usar \texttt{shapes.geometric} y flechas \texttt{Stealth}.

\subsubsection*{Diagrama de Venn}

Muestra coincidencias y diferencias entre conjuntos. Funciona para comparar derecho/moral, legalidad/justicia, reglas/principios o corrientes filosóficas.

\textbf{Construcción:} definir conjuntos, colocar rasgos exclusivos y zona de intersección. En TikZ usar círculos semitransparentes.

\subsubsection*{Diagrama de Gowin}

Relaciona pregunta central, conceptos, principios, registros, transformaciones y conclusiones. Es útil para analizar un problema de investigación o una norma desde teoría y evidencia.

\textbf{Construcción:} formular pregunta, colocar teoría de un lado y evidencia del otro, derivar conclusión. En TikZ se puede construir con una V central y bloques laterales.

\subsubsection*{Mapa de cajas, mapa semántico y redes conceptuales}

Estos productos organizan conceptos por agrupaciones, campos de significado o conexiones. Son útiles para conceptos jurídicos fundamentales o familias de teorías.

\textbf{Construcción:} agrupar por criterio, no por ocurrencia. Cada caja o nodo debe tener función. En TikZ usar cajas por bloque, colores moderados y flechas solo cuando exista relación real.

\subsubsection*{Mapeo de procesos}

Sirve para describir un sistema operativo: entradas, actividades, responsables, controles y salidas. En Derecho puede aplicarse a etapas de interpretación, aplicación normativa o solución de controversias.

\textbf{Construcción:} definir entrada, proceso, control, salida y retroalimentación. En TikZ usar flujo horizontal o carriles.

\subsubsection*{Gráfico estadístico e histograma}

Usar cuando existan datos cuantitativos. No inventar cifras. Si no hay datos confiables, no usar gráfico estadístico como decoración.

\textbf{Construcción:} definir variable, fuente, escala, unidades y lectura del dato. En LaTeX puede usarse \texttt{tikzpicture} o \texttt{pgfplots} si se habilita el paquete.

\subsubsection*{SQA}

SQA organiza tres momentos: qué sé, qué quiero saber y qué aprendí. Es útil para iniciar y cerrar una investigación.

\textbf{Construcción:} usar tabla de tres columnas. La última columna debe mostrar aprendizaje nuevo, no repetir la primera.

\subsubsection*{Valoración de decisiones}

Permite comparar alternativas con criterios cualitativos o cuantitativos. Es útil para justificar una postura entre corrientes, métodos interpretativos o soluciones de caso.

\textbf{Construcción:} definir alternativas, establecer criterios, asignar pesos, valorar y justificar la decisión. En LaTeX usar tabla o matriz ponderada.

\subsubsection*{Ensayo, informe, artículo, monografía y reporte de investigación}

Son productos escritos. No requieren TikZ salvo que incorporen figuras. Deben tener estructura: introducción con problema, desarrollo con fuentes y análisis, conclusión con postura, y referencias.

\textbf{Construcción:} planteamiento, marco conceptual, análisis, aplicación, postura y cierre. Evitar texto enciclopédico. La voz debe convertir información en criterio.

\subsubsection*{Reseña, resumen, reporte de lectura, sumillado y exegética}

Son productos de lectura y síntesis. Deben demostrar comprensión del texto fuente, no sustituirlo por paráfrasis superficial.

\textbf{Construcción:} identificar tesis, argumentos, conceptos clave, límite del texto y lectura personal. Citar correctamente.

\subsubsection*{Debate, foro, panel, mesa redonda, controversia y suasoria}

Son productos argumentativos. Requieren postura, contraargumentos y cierre. En Derecho son útiles para defender interpretaciones o criterios.

\textbf{Construcción:} tesis, argumentos, evidencia, objeción, respuesta y conclusión. Si se presenta por escrito, usar subtítulos claros.

\subsubsection*{Estudio de caso, análisis de hechos y noticia falsa}

Estos productos trabajan con hechos. Deben distinguir hecho, interpretación, norma/principio, problema y solución.

\textbf{Construcción:} hechos relevantes, problema central, criterios de análisis, aplicación, postura y consecuencia.

\subsubsection*{Cartel, afiche, anuncio, fotomontaje y videocápsula}

Son productos comunicativos y visuales. Si la entrega es PDF, el cartel o afiche puede construirse con TikZ. Si se requiere video, el documento debe incluir guion, enlace, captura y justificación.

\textbf{Construcción:} objetivo, público, mensaje central, jerarquía visual, fuentes y cierre. No saturar texto.

\subsubsection*{Glosario colaborativo}

Debe incluir término, definición propia, fuente, ejemplo y comentario. No basta copiar definiciones.

\textbf{Construcción:} ordenar alfabéticamente o por familias conceptuales; citar fuentes; agregar aplicación jurídica.

\subsubsection*{Portafolio de evidencias digital}

Integra productos y reflexión sobre el proceso. Debe mostrar trayectoria, no solo acumulación.

\textbf{Construcción:} índice, evidencias, explicación de cada evidencia, aprendizaje, correcciones y conclusión.

\subsubsection*{Simulación, sociodrama, juego de roles, juego de negocios y taller}

Son técnicas de actuación o aplicación. Si se entregan por escrito, deben incluir guion, roles, contexto, desarrollo, resultados y reflexión.

\textbf{Construcción:} problema, roles, reglas, escena o dinámica, decisiones, cierre y aprendizaje.

\subsubsection*{Consulta pública, encuesta, grupos focales y entrevista}

Son técnicas de recolección de información. Requieren objetivo, participantes, instrumento, resultados y análisis.

\textbf{Construcción:} delimitar problema, formular preguntas claras, aplicar, organizar resultados, interpretar y cerrar con hallazgos.

\subsection{Catálogo de las 100 técnicas}

\begin{longtable}{@{}p{0.08\linewidth}p{0.35\linewidth}p{0.16\linewidth}p{0.31\linewidth}@{}}
\toprule
\textbf{No.} & \textbf{Técnica} & \textbf{Nivel} & \textbf{Producto sugerido en LaTeX/TikZ} \\
\midrule
\endhead
1 & Cuadro comparativo & Sintetizar & Longtable o matriz TikZ. \\
2 & Cuadro sinóptico & Sintetizar & Árbol jerárquico TikZ. \\
3 & Cuestionario & Explicar & Tabla de preguntas/respuestas. \\
4 & Debate & Aplicar & Guion argumentativo. \\
5 & Diagrama de Gowin & Aplicar & Diagrama V en TikZ. \\
6 & Ejemplificación & Explicar & Tabla concepto-ejemplo-aplicación. \\
7 & Entrevista & Aplicar & Guion y matriz de hallazgos. \\
8 & Esquema & Sintetizar & Cajas jerárquicas TikZ. \\
9 & Glosario colaborativo & Explicar & Tabla término-definición-fuente. \\
10 & Historieta & Recordar & Viñetas TikZ o guion. \\
11 & Informe & Sintetizar & Documento con secciones. \\
12 & Línea de tiempo & Sintetizar & Timeline TikZ. \\
13 & Mapa conceptual & Sintetizar & Nodos y conectores TikZ. \\
14 & Panel de discusión & Aplicar & Tabla de posturas. \\
15 & Pensamiento de diseño & Recordar & Fases empatizar-definir-idear-prototipar-evaluar. \\
16 & Phillips 66 & Aplicar & Matriz grupo-idea-síntesis. \\
17 & Resumen & Recordar & Texto breve estructurado. \\
18 & Simulación & Construir & Escenario, variables y resultados. \\
19 & Sociodrama & Aplicar & Guion de escenas. \\
20 & Trabajo cooperativo & Aplicar & Plan de roles y evidencias. \\
21 & Análisis de contenido & Analizar & Matriz texto-idea-función-crítica. \\
22 & Análisis de tergiversación textual & Analizar & Tabla afirmación/evidencia/corrección. \\
23 & Análisis y consensos & Analizar & Matriz de acuerdos y disensos. \\
24 & Autoexplicación & Explicar & Secuencia pregunta-respuesta-criterio. \\
25 & Cronología ilustrada & Construir & Línea de tiempo con iconos. \\
26 & Estructuras textuales & Analizar & Mapa de estructura argumentativa. \\
27 & Estudio de casos & Explicar & Hechos-problema-norma-solución. \\
28 & Exposición oral & Aplicar & Guion o diapositivas. \\
29 & Foro & Aplicar & Entrada, réplica y cierre. \\
30 & Grupos de discusión & Aplicar & Matriz de aportaciones. \\
31 & Grupos focales & Explicar & Guion y síntesis de hallazgos. \\
32 & Mapa de cajas & Sintetizar & Cajas agrupadas TikZ. \\
33 & Mapa mental & Explicar & Nodos radiales TikZ. \\
34 & Monografía & Explicar & Documento de investigación. \\
35 & Preguntas y premios & Recordar & Banco de reactivos. \\
36 & Reporte de lectura general & Explicar & Ficha de lectura ampliada. \\
37 & Reseña & Sintetizar & Valoración crítica breve. \\
38 & Rompecabezas & Construir & Integración por piezas. \\
39 & Sumillado & Sintetizar & Texto con notas marginales. \\
40 & Testimonio & Recordar & Entrevista y relato analizado. \\
41 & Consulta pública & Aplicar & Instrumento y matriz de resultados. \\
42 & Controversia estructurada & Analizar & Tesis, antítesis y síntesis. \\
43 & Debate público & Aplicar & Guion de argumentación. \\
44 & Diagrama de Venn & Analizar & Círculos TikZ. \\
45 & Diario de aprendizaje & Recordar & Bitácora reflexiva. \\
46 & Discusión de gabinete & Analizar & Tabla de roles y decisiones. \\
47 & Documentales & Explicar & Guion audiovisual. \\
48 & Encuesta interactiva & Aplicar & Instrumento y gráfica. \\
49 & Ensayo & Construir & Texto argumentativo. \\
50 & Estudio de noticia falsa & Analizar & Verificación y contraste de fuentes. \\
51 & Exegética & Recordar & Interpretación textual guiada. \\
52 & Fábula & Recordar & Narración con moraleja. \\
53 & Heurística de Bruner & Analizar & Ruta descubrimiento-concepto. \\
54 & Histograma & Analizar & Gráfico de frecuencias. \\
55 & Pensamiento analógico & Aplicar & Analogía y transferencia. \\
56 & Poesía lírica & Construir & Producto creativo con reflexión. \\
57 & Sesión bibliográfica & Analizar & Fichas por fuente. \\
58 & SQA & Recordar & Tabla sé-quiero-aprendí. \\
59 & Suasoria & Aplicar & Discurso persuasivo. \\
60 & Tablón de anuncios & Sintetizar & Muro informativo o tablero. \\
61 & Afiche & Explicar & Cartel breve en TikZ. \\
62 & Análisis de hechos & Analizar & Hechos-causas-consecuencias. \\
63 & Anuncio publicitario & Explicar & Pieza persuasiva. \\
64 & Atributos & Explicar & Tabla de cualidades. \\
65 & Diagrama de flujo & Sintetizar & Flujo TikZ. \\
66 & Diálogos simultáneos & Aplicar & Guion de interacción. \\
67 & Feynman & Explicar & Explicación simple y prueba. \\
68 & Gráfico estadístico & Analizar & Gráfico con fuente. \\
69 & Lluvia de ideas dirigida & Explicar & Mapa radial o lista clasificada. \\
70 & Mapa semántico & Sintetizar & Red de significados TikZ. \\
71 & Mapeo de procesos & Analizar & Proceso entrada-salida. \\
72 & Mesa redonda con interrogador & Aplicar & Matriz de preguntas/posturas. \\
73 & Mnemotecnia & Recordar & Reglas de memoria. \\
74 & Portafolio de evidencias digital & Construir & Índice de evidencias. \\
75 & Redes conceptuales & Sintetizar & Grafo TikZ. \\
76 & Reporte de investigación & Explicar & Documento con método. \\
77 & Simposio & Aplicar & Programa y síntesis. \\
78 & Socioaprendizaje & Aplicar & Comunidad y evidencias. \\
79 & Tuits & Sintetizar & Microargumentos. \\
80 & Valoración de decisiones & Analizar & Matriz ponderada. \\
81 & Artículo & Analizar & Texto con tesis y fuentes. \\
82 & Cartel & Explicar & Cartel TikZ. \\
83 & Demostración silenciosa & Aplicar & Secuencia visual. \\
84 & Descripción de un personaje & Aplicar & Perfil estructurado. \\
85 & Estudio intercalado & Aplicar & Plan de estudio espaciado. \\
86 & Examen práctico & Aplicar & Evidencia de ejecución. \\
87 & Fotomontaje digital & Construir & Composición visual justificada. \\
88 & Instrucción personalizada & Analizar & Ruta de aprendizaje individual. \\
89 & Journal digital & Construir & Bitácora digital. \\
90 & Juego de negocios & Analizar & Simulación estratégica. \\
91 & Juego de roles & Aplicar & Guion de roles. \\
92 & Práctica distribuida & Aplicar & Calendario de práctica. \\
93 & Preguntas dirigidas & Construir & Banco de preguntas guía. \\
94 & Proyectos colaborativos & Construir & Plan de proyecto. \\
95 & Reportaje & Construir & Investigación narrativa. \\
96 & Scamper & Recordar & Matriz sustituir-combinar-adaptar. \\
97 & Social media & Aplicar & Estrategia de publicación. \\
98 & Taller & Aplicar & Secuencia de actividades. \\
99 & Videocápsula & Sintetizar & Guion y storyboard. \\
100 & Visitas guiadas virtuales 360 & Construir & Guion, secuencia y recorrido. \\
\bottomrule
\end{longtable}

\subsection{Plantillas TikZ mínimas}

Estas plantillas se pueden copiar a la actividad concreta y ajustar.

\subsubsection*{Mapa conceptual}

\begin{verbatim}
\begin{figure}[H]
\centering
\begin{tikzpicture}[
  concepto/.style={rectangle, draw, rounded corners=2pt,
    align=center, text width=3.2cm, minimum height=0.9cm},
  flecha/.style={-{Stealth[length=2mm]}, thick}
]
\node[concepto] (a) {Concepto central};
\node[concepto, below left=1.3cm and 1.8cm of a] (b) {Concepto A};
\node[concepto, below right=1.3cm and 1.8cm of a] (c) {Concepto B};
\draw[flecha] (a) -- node[above, sloped]{fundamenta} (b);
\draw[flecha] (a) -- node[above, sloped]{se relaciona con} (c);
\end{tikzpicture}
\caption{Mapa conceptual de [tema]}
\end{figure}
\end{verbatim}

\subsubsection*{Línea de tiempo}

\begin{verbatim}
\begin{figure}[H]
\centering
\begin{tikzpicture}[
  evento/.style={rectangle, draw, rounded corners=2pt,
    align=center, text width=3cm, font=\scriptsize},
  eje/.style={-{Stealth[length=2mm]}, thick}
]
\draw[eje] (0,0) -- (12,0);
\foreach \x/\anio/\texto in {
  0/1917/Constitucion,
  4/2011/Reforma DDHH,
  8/2013/Ley de Amparo,
  12/2024/Reforma vigente}
{
  \draw (\x,0.12) -- (\x,-0.12);
  \node[evento, above=0.45cm] at (\x,0) {\anio\\\texto};
}
\end{tikzpicture}
\caption{Linea de tiempo de [tema]}
\end{figure}
\end{verbatim}

\subsubsection*{Diagrama de flujo}

\begin{verbatim}
\begin{figure}[H]
\centering
\begin{tikzpicture}[
  paso/.style={rectangle, draw, rounded corners=2pt,
    align=center, text width=3.1cm, minimum height=0.9cm},
  decision/.style={diamond, draw, aspect=2, align=center,
    text width=2.8cm, inner sep=1pt},
  flecha/.style={-{Stealth[length=2mm]}, thick}
]
\node[paso] (inicio) {Identificar problema};
\node[paso, right=1.4cm of inicio] (norma) {Localizar norma};
\node[decision, right=1.4cm of norma] (decision) {Hay conflicto?};
\node[paso, below=1.2cm of decision] (analisis) {Argumentar solucion};
\draw[flecha] (inicio) -- (norma);
\draw[flecha] (norma) -- (decision);
\draw[flecha] (decision) -- node[right]{si} (analisis);
\end{tikzpicture}
\caption{Diagrama de flujo de [proceso]}
\end{figure}
\end{verbatim}

\subsection{Cierre operativo}

La elección de una técnica no debe ser decorativa. Si la actividad pide un producto, la técnica funciona como una herramienta de análisis. El criterio de calidad es que el producto permita ver una decisión intelectual: qué se seleccionó, cómo se organizó, qué relación se construyó y qué postura se puede defender a partir de esa organización.
